% =============================================================================
%                       Chapter 4: Software Development
% =============================================================================



\setcounter{chapter}{4}
\setcounter{figure}{0}
\setcounter{table}{0}
\chapter*{Chapter 4}
\markboth{Software Development}{Software Development} 
\section{Software Development}

This chapter describes the implementation of the Smart AI Powered 
Agriculture System at a finer level of detail. It explains the 
development standards, environment, and key software modules across 
the embedded firmware (ESP32), backend services (Node.js), mobile 
application (Flutter), and the AI/ML recommendation service 
\cite{thilakarathne2023fields}. The implementation is aligned with 
the system design described in Chapter 3.

% ----------------------------------------------------------
% 4.1 Coding Standards
% ----------------------------------------------------------
\subsection{Coding Standards}
\label{sec:coding-standards}

To ensure maintainability, readability, and consistent collaboration 
among team members, the following coding standards were followed 
\cite{sommerville2016software}:

\subsubsection{General Standards}
Consistent naming conventions were used across all modules (camelCase for variables/functions, PascalCase for classes). Functions were kept small and single-purpose to improve modularity and testability \cite{sommerville2016software}. All modules include meaningful comments for complex logic and configuration parameters. Sensitive information (API keys, tokens) was kept out of source code and stored using environment variables or secure configuration files \cite{jwtauth2017}. Proper error handling and validation were applied at the boundaries (API inputs, sensor readings, and database writes).


\subsubsection{ESP32 Firmware Standards}
Clear separation of concerns: sensor reading, networking, and actuation logic implemented as separate functions \cite{hassebo2024esp32}. Debouncing and filtering used for noisy sensors (e.g., rain sensor) to avoid false triggers \cite{noisyiot2023}. Safe defaults: irrigation is turned \textbf{OFF} when network or sensor failures occur to avoid uncontrolled watering \cite{hassebo2024esp32}.


\subsubsection{Backend Standards (Node.js)}
REST endpoints follow consistent URI patterns and HTTP status codes \cite{thilakarathne2023fields}. Middleware-based validation and authentication (JWT) is enforced for protected routes \cite{jwtauth2017}. Logging includes request IDs and timestamps for debugging and audit.


\subsubsection{Mobile Application Standards}
UI components follow a consistent design system for readability and ease of use \cite{mobileui2022}. Data is fetched asynchronously and cached for better user experience. Input forms (thresholds, schedules) validate values before sending to backend.


% ----------------------------------------------------------
% 4.2 Development Environment
% ----------------------------------------------------------
\subsection{Development Environment}
\label{sec:dev-environment}

This section describes the tools, platforms, and configurations 
used for software development.

\subsubsection{Tools and Platforms}
\textbf{Embedded/Firmware:} Arduino IDE with ESP32 board support packages \cite{raj2021survey}. \textbf{Backend:} Node.js runtime with Express.js framework \cite{thilakarathne2023fields}. \textbf{Database:} MySQL (MariaDB) for all data storage and retrieval operations \cite{uotila2020db}. \textbf{Mobile:} Flutter SDK (cross-platform) \cite{khawas2018firebase}. \textbf{AI/ML:} Python environment Random Forest \cite{ali2025soilml}. \textbf{Version Control:} Git for collaborative development and change tracking \cite{sommerville2016software}.


\subsubsection{Testing Setup}
\textbf{Firmware testing:} Serial monitor logging and controlled sensor value simulation \cite{hassebo2024esp32}. \textbf{Backend testing:} API testing using Postman and unit checks for validation logic. \textbf{Mobile testing:} Emulators and physical device testing for UI and notifications \cite{khawas2018firebase}. \textbf{End-to-end testing:} Real sensor readings sent to backend, verified on app dashboard and alert system \cite{emon2025iot}.


% ----------------------------------------------------------
% 4.3 Software Description
% ----------------------------------------------------------
\subsection{Software Description}
\label{sec:software-description}

The Smart AI Powered Agriculture System is fully functional and 
deployed. The core implementation is based on ESP32 firmware that 
integrates sensors, Wi-Fi communication, backend API interaction, 
local safety control, and relay-based irrigation automation 
\cite{hassebo2024esp32}. The firmware is responsible for reading 
environmental data, applying local decision logic, receiving pump 
commands from the backend, and uploading sensor readings for 
storage, alerts, analytics, and AI-based processing 
\cite{thilakarathne2023fields}.

\vspace{0.5em}
\noindent\textbf{Note on Code Snippets:} The following code snippets are presented to highlight the core functional logic (e.g., sensor integration, decision-making, and communication protocols). For the purpose of brevity and readability, standard imports, module definitions, and boilerplate file structures have been intentionally omitted. For access to the complete, fully functional source code across all system layers, please refer to the project's official GitHub repository \cite{smart_agri_github}.

% ==========================================================
% 4.3.1 Snippet 1 — Hardware Pin Configuration
% ==========================================================
\subsubsection{Snippet 1: Hardware Pin Configuration}
\label{subsec:snippet1}

As implemented in \texttt{ESP32\_Firmware/final ESP32 code.txt} (Lines 34--44), the ESP32 firmware defines dedicated GPIO pins for soil moisture, 
rainfall, light intensity, temperature/humidity, and relay-based 
pump control. This explicit pin mapping connects the physical field unit 
with the software logic \cite{raj2021survey}.

\begin{lstlisting}
// Pin Definitions
#define SOIL_PIN  34
#define RAIN_PIN  33
#define LIGHT_PIN 32
#define DHT_PIN   4
#define RELAY_PIN 5

#define DHT_TYPE DHT22
DHT dht(DHT_PIN, DHT_TYPE);
\end{lstlisting}

\textbf{Outcome:} The hardware components are mapped with ESP32 
pins, enabling the firmware to collect sensor data and control the 
irrigation relay.

% ==========================================================
% 4.3.2 Snippet 2 — Sensor Threshold Configuration
% ==========================================================
\subsubsection{Snippet 2: Sensor Threshold Configuration}
\label{subsec:snippet2}

Referencing \texttt{ESP32\_Firmware/final ESP32 code.txt} (Lines 46--52), hardcoded threshold values are defined to classify soil conditions, rainfall 
status, and light intensity. These limits are utilized by the firmware 
for local offline automation and zero-latency safety decisions.

\begin{lstlisting}
// Sensor Thresholds
#define SOIL_DRY_THRESHOLD   2800
#define SOIL_WET_THRESHOLD   1200
#define RAIN_THRESHOLD       2000
#define LIGHT_DARK_THRESHOLD 3000
\end{lstlisting}

\textbf{Outcome:} The system can identify dry soil, wet soil, 
rainfall, and dark light conditions using predefined sensor 
thresholds.

% ==========================================================
% 4.3.3 Snippet 3 — Three-Tier Timing System
% ==========================================================
\subsubsection{Snippet 3: Three-Tier Timing System}
\label{subsec:snippet3}

As shown in \texttt{ESP32\_Firmware/final ESP32 code.txt} (Lines 73--78), the firmware establishes a strict three-tier timing mechanism. This architecture separates fast local safety checks (1s), backend command polling (5s), and full telemetry data uploads (30s) to optimize network bandwidth \cite{hassebo2024esp32}.

\begin{lstlisting}
// Timing Configuration
#define SENSOR_READ_INTERVAL   1000
#define COMMAND_POLL_INTERVAL  5000
#define SENSOR_UPLOAD_INTERVAL 30000
#define BACKEND_RETRY_INTERVAL 5000
#define WIFI_RETRY_INTERVAL_MS 10000
#define PUMP_MIN_RUN_TIME      1500
\end{lstlisting}

\textbf{Outcome:} Sensors are read every second, backend commands 
are checked every five seconds, and complete sensor data is 
uploaded every thirty seconds.

% ==========================================================
% 4.3.4 Snippet 4 — Safe System Initialization
% ==========================================================
\subsubsection{Snippet 4: Safe System Initialization}
\label{subsec:snippet4}

Located in \texttt{ESP32\_Firmware/final ESP32 code.txt} (Lines 131--168), the system startup sequence initializes serial communication, starts 
the DHT sensor, configures the relay pin, and explicitly keeps the pump OFF 
by default. Since the relay operates via an active-low trigger, HIGH safely enforces the OFF state \cite{hassebo2024esp32}.

\begin{lstlisting}
void setup() {
  Serial.begin(115200);
  delay(2000);

  dht.begin();

  // Active-LOW relay: HIGH = OFF
  pinMode(RELAY_PIN, OUTPUT_OPEN_DRAIN);
  digitalWrite(RELAY_PIN, HIGH);

  connectToWiFi();
}
\end{lstlisting}

\textbf{Outcome:} The pump remains safely OFF during startup, 
preventing accidental irrigation when the device powers on.

% ==========================================================
% 4.3.5 Snippet 5 — Main Execution Loop
% ==========================================================
\subsubsection{Snippet 5: Main Execution Loop}
\label{subsec:snippet5}

The main execution flow, implemented in \texttt{ESP32\_Firmware/final ESP32 code.txt} (Lines 173--230), drives the continuous monitoring loop. It dynamically checks 
Wi-Fi status, reads sensors, evaluates local safety logic, polls 
backend endpoints for irrigation commands, and handles telemetry uploads \cite{emon2025iot}.

\begin{lstlisting}
void loop() {
  unsigned long now = millis();

  if (WiFi.status() == WL_CONNECTED) {
    wifiConnected = true;
  } else {
    wifiConnected = false;
    if (now - lastWifiAttemptMs >= WIFI_RETRY_INTERVAL_MS) {
      connectToWiFi();
      lastWifiAttemptMs = millis();
    }
  }

  if (now - lastSensorReadTime >= SENSOR_READ_INTERVAL) {
    lastSensorReadTime = now;
    readAllSensors(cachedSoilValue, cachedRainValue,
                   cachedLightValue, cachedTemp, cachedHum);

    if (wifiConnected)
      applyLocalRainSafety(cachedRainValue);
    else
      runOfflinePumpLogic(cachedSoilValue, cachedRainValue);
  }

  if (wifiConnected &&
      now - lastCommandPollTime >= COMMAND_POLL_INTERVAL) {
    lastCommandPollTime = now;
    fetchPumpCommand();
  }

  if (wifiConnected &&
      now - lastUploadTime >= SENSOR_UPLOAD_INTERVAL) {
    if (sendSensorDataToBackend(cachedSoilValue, cachedRainValue,
        cachedLightValue, cachedTemp, cachedHum)) {
      lastUploadTime = now;
    }
  }
}
\end{lstlisting}

\textbf{Outcome:} The firmware performs real-time monitoring, 
backend synchronization, and sensor data upload in a structured 
and efficient manner.

% ==========================================================
% 4.3.6 Snippet 6 — Smoothed Sensor Reading
% ==========================================================
\subsubsection{Snippet 6: Smoothed Sensor Reading}
\label{subsec:snippet6}

To mitigate sensor noise, the firmware applies a smoothing algorithm found in \texttt{ESP32\_Firmware/final ESP32 code.txt} (Lines 346--353). Analog sensor values are averaged across multiple sequential samples to improve reading stability and prevent false condition triggers \cite{noisyiot2023}.

\begin{lstlisting}
int readAnalogSmoothed(uint8_t pin) {
  long sum = 0;

  for (int i = 0; i < ANALOG_SAMPLES; i++) {
    sum += analogRead(pin);
    delayMicroseconds(150);
  }

  return (int)(sum / ANALOG_SAMPLES);
}
\end{lstlisting}

\textbf{Outcome:} Sensor readings become more stable and reliable, 
reducing false triggers caused by noisy analog signals.

% ==========================================================
% 4.3.7 Snippet 7 — Sensor Data Acquisition
% ==========================================================
\subsubsection{Snippet 7: Sensor Data Acquisition}
\label{subsec:snippet7}

As implemented in \texttt{ESP32\_Firmware/final ESP32 code.txt} (Lines 355--371), the firmware acquires soil moisture, rainfall, light intensity, 
temperature, and humidity metrics from the connected hardware interfaces 
\cite{mohanraj2016field}.

\begin{lstlisting}
bool readDht(float& tempC, float& humPct) {
  humPct = dht.readHumidity();
  tempC  = dht.readTemperature();

  if (isnan(humPct) || isnan(tempC)) {
    tempC  = -999;
    humPct = -999;
    return false;
  }

  return true;
}

void readAllSensors(int& soil, int& rain, int& light,
                    float& tempC, float& humPct) {
  soil  = readAnalogSmoothed(SOIL_PIN);
  rain  = readAnalogSmoothed(RAIN_PIN);
  light = readAnalogSmoothed(LIGHT_PIN);
  readDht(tempC, humPct);
}
\end{lstlisting}

\textbf{Outcome:} The ESP32 collects all required environmental 
readings and stores them in cached variables for processing and 
upload.

% ==========================================================
% 4.3.8 Snippet 8 — Local Rain Safety Mechanism
% ==========================================================
\subsubsection{Snippet 8: Local Rain Safety Mechanism}
\label{subsec:snippet8}

Detailed in \texttt{ESP32\_Firmware/final ESP32 code.txt} (Lines 238--263), the system incorporates a zero-latency local rain override. If rainfall 
is detected, the irrigation pump is immediately forced OFF without requiring network 
communication to the backend API \cite{garcia2020iot}.

\begin{lstlisting}
void applyLocalRainSafety(int rainValue) {
  bool raining = (rainValue < RAIN_THRESHOLD);

  if (raining) {
    if (!rainOverrideActive) {
      rainOverrideActive = true;
      turnPumpOff();
      Serial.println("Rain detected: pump OFF");
    }
    return;
  }

  if (rainOverrideActive) {
    rainOverrideActive = false;
    int restoreCmd = lastBackendPumpState;
    lastBackendPumpState = -2;
    syncRelayToCommand(restoreCmd, "rain_cleared_restore");
  }
}
\end{lstlisting}

\textbf{Outcome:} The irrigation pump is protected from running 
during rainfall, reducing water wastage and improving system safety.

% ==========================================================
% 4.3.9 Snippet 9 — Backend Command Polling
% ==========================================================
\subsubsection{Snippet 9: Backend Command Polling}
\label{subsec:snippet9}

As coded in \texttt{ESP32\_Firmware/final ESP32 code.txt} (Lines 270--312), the ESP32 polls the backend via an HTTP GET request to retrieve the latest pump command sent 
from the mobile application or backend automation logic \cite{yokotani2016mqtt}.

\begin{lstlisting}
bool fetchPumpCommand() {
  if (WiFi.status() != WL_CONNECTED) return false;

  WiFiClientSecure secureClient;
  HTTPClient http;

  String endpoint = String("https://") + BACKEND_HOST +
                    COMMAND_PATH + deviceId;

  secureClient.setInsecure();
  if (!http.begin(secureClient, endpoint)) return false;

  int code = http.GET();

  if (code == 200) {
    String response = http.getString();
    StaticJsonDocument<128> doc;

    if (!deserializeJson(doc, response) &&
        doc.containsKey("pump_status")) {
      int commanded = doc["pump_status"].as<int>();
      const char* reason = doc["pump_reason"] | "unknown";
      syncRelayToCommand(commanded, reason);
    }
  }

  http.end();
  return code == 200;
}
\end{lstlisting}

\textbf{Outcome:} The field device stays synchronized with backend 
pump commands and supports remote irrigation control from the 
application.

% ==========================================================
% 4.3.10 Snippet 10 — Relay Synchronization Logic
% ==========================================================
\subsubsection{Snippet 10: Relay Synchronization Logic}
\label{subsec:snippet10}

Located in \texttt{ESP32\_Firmware/final ESP32 code.txt} (Lines 320--341), the relay synchronization function applies backend commands only 
when the desired pump state differs from the current state. Furthermore, it explicitly blocks remote ON commands if the local rain override is active \cite{hassebo2024esp32}.

\begin{lstlisting}
void syncRelayToCommand(int commanded, const char* reason) {
  if (rainOverrideActive && commanded == 1) {
    Serial.println("Command blocked due to rain override");
    return;
  }

  if (commanded == lastBackendPumpState) {
    return;
  }

  lastBackendPumpState = commanded;

  if (commanded == 1 && !pumpRunning) {
    turnPumpOn();
  } else if (commanded == 0 && pumpRunning) {
    turnPumpOff();
  }
}
\end{lstlisting}

\textbf{Outcome:} The relay avoids unnecessary repeated switching 
and ensures backend commands cannot override local rain safety.

% ==========================================================
% 4.3.11 Snippet 11 — Relay-Based Pump Control
% ==========================================================
\subsubsection{Snippet 11: Relay-Based Pump Control}
\label{subsec:snippet11}

Described in \texttt{ESP32\_Firmware/final ESP32 code.txt} (Lines 463--472), the irrigation pump is controlled via an active-low relay module. Setting the 
relay pin LOW turns the pump ON, while setting it HIGH turns the pump OFF.

\begin{lstlisting}
void turnPumpOn() {
  digitalWrite(RELAY_PIN, LOW);
  pumpRunning = true;
  pumpStartTime = millis();
}

void turnPumpOff() {
  digitalWrite(RELAY_PIN, HIGH);
  pumpRunning = false;
}
\end{lstlisting}

\textbf{Outcome:} The system can physically control the irrigation 
pump through the ESP32 and relay module.

% ==========================================================
% 4.3.12 Snippet 12 — Offline Irrigation Logic
% ==========================================================
\subsubsection{Snippet 12: Offline Irrigation Logic}
\label{subsec:snippet12}

Defined in \texttt{ESP32\_Firmware/final ESP32 code.txt} (Lines 412--458), an offline fallback mechanism ensures that if Wi-Fi is disconnected, the ESP32 autonomously continues irrigation control 
locally using its direct soil moisture and rainfall sensor inputs 
\cite{emon2025iot}.

\begin{lstlisting}
void runOfflinePumpLogic(int soilValue, int rainValue) {
  bool raining = (rainValue < RAIN_THRESHOLD);
  bool soilWet = (soilValue < SOIL_WET_THRESHOLD);
  bool soilDry = (soilValue > SOIL_DRY_THRESHOLD);

  if (soilDry) requestedPumpState = true;
  else if (soilWet) requestedPumpState = false;

  if (raining) {
    rainOverrideActive = true;
    if (pumpRunning) turnPumpOff();
    return;
  }

  if (soilWet && pumpRunning &&
      millis() - pumpStartTime >= PUMP_MIN_RUN_TIME) {
    turnPumpOff();
    return;
  }

  if (requestedPumpState && !pumpRunning) {
    turnPumpOn();
  }
}
\end{lstlisting}

\textbf{Outcome:} The irrigation system remains functional even 
when internet connectivity is unavailable.

% ==========================================================
% 4.3.13 Snippet 13 — Wi-Fi Connectivity and Reconnection
% ==========================================================
\subsubsection{Snippet 13: Wi-Fi Connectivity and Reconnection}
\label{subsec:snippet13}

As seen in \texttt{ESP32\_Firmware/final ESP32 code.txt} (Lines 477--495), the ESP32 connects to Wi-Fi in station mode and updates its global connection state flag. If the connection drops or fails, the firmware bypasses network uploads and routes execution to the offline safety logic \cite{farhan2019iotplatform}.

\begin{lstlisting}
void connectToWiFi() {
  if (WiFi.status() == WL_CONNECTED) {
    wifiConnected = true;
    return;
  }

  WiFi.mode(WIFI_STA);
  WiFi.begin("Your_SSID", "Your_Password");

  unsigned long startTime = millis();
  while (WiFi.status() != WL_CONNECTED &&
         millis() - startTime < 12000) {
    delay(400);
  }

  wifiConnected = (WiFi.status() == WL_CONNECTED);
}
\end{lstlisting}

\textbf{Outcome:} The firmware supports automatic Wi-Fi connection 
handling while maintaining system operation in offline mode.

% ==========================================================
% 4.3.14 Snippet 14 — Sensor Value Conversion
% ==========================================================
\subsubsection{Snippet 14: Sensor Value Conversion}
\label{subsec:snippet14}

Detailed in \texttt{ESP32\_Firmware/final ESP32 code.txt} (Lines 500--506), raw 12-bit analog values (0--4095) are linearly mapped and constrained into 
percentage values (0--100\%) before transmission to the backend, standardizing the data format.

\begin{lstlisting}
float convertSoilMoistureToPercentage(int rawValue) {
  return constrain((float)map(rawValue, 0, 4095, 100, 0),
                   0.0f, 100.0f);
}

float convertLightToPercentage(int rawValue) {
  return constrain((float)map(rawValue, 0, 4095, 100, 0),
                   0.0f, 100.0f);
}
\end{lstlisting}

\textbf{Outcome:} Raw ADC readings are normalized into 
user-friendly percentage values for dashboard visualization and 
backend processing.

% ==========================================================
% 4.3.15 Snippet 15 — Sensor Data Upload to Backend
% ==========================================================
\subsubsection{Snippet 15: Sensor Data Upload to Backend}
\label{subsec:snippet15}

Finally, in \texttt{ESP32\_Firmware/final ESP32 code.txt} (Lines 530--605), the firmware constructs a JSON payload containing the device ID, normalized sensor 
values, rainfall status, and current pump state. This telemetry data is uploaded to 
the Node.js backend API using an HTTP POST request \cite{yokotani2016mqtt}.

\begin{lstlisting}
bool sendSensorDataToBackend(int soilRaw, int rainRaw,
                             int lightRaw, float temp, 
                             float hum) {
  if (WiFi.status() != WL_CONNECTED) return false;

  float soilMoisture = 
    convertSoilMoistureToPercentage(soilRaw);
  float lightIntensity = 
    convertLightToPercentage(lightRaw);
  bool raining = (rainRaw < RAIN_THRESHOLD);

  StaticJsonDocument<384> doc;
  doc["device_id"]      = deviceId;
  doc["soil_moisture"]  = soilMoisture;
  doc["light_intensity"]= lightIntensity;
  doc["rainfall"]       = raining;
  doc["pump_on"]        = pumpRunning ? 1 : 0;

  if (temp != -999) doc["temperature"] = temp;
  if (hum  != -999) doc["humidity"]    = hum;

  String jsonString;
  serializeJson(doc, jsonString);

  WiFiClientSecure secureClient;
  HTTPClient http;

  String endpoint = String("https://") + 
                    BACKEND_HOST + API_PATH;
  secureClient.setInsecure();

  http.begin(secureClient, endpoint);
  http.addHeader("Content-Type", "application/json");

  int code = http.POST(jsonString);
  http.end();

  return (code == 200 || code == 201);
}
\end{lstlisting}

\textbf{Outcome:} Sensor readings are successfully transmitted to 
the backend for database storage, alerts, analytics, and 
application display.

% ----------------------------------------------------------
% 4.4 Implementation Challenges and Resolutions
% ----------------------------------------------------------
\subsection{Implementation Challenges and Resolutions}
\label{sec:challenges}

During implementation, several challenges were encountered and 
resolved \cite{noisyiot2023}:
\textbf{Noisy Sensor Readings:} Soil moisture and rain sensors produced unstable values in early tests. This was resolved by applying filtering, threshold smoothing, and debouncing logic \cite{noisyiot2023}. \textbf{Intermittent Wi-Fi Connectivity:} Network dropouts caused missed telemetry uploads. Automatic reconnect and retry mechanisms were introduced in firmware \cite{farhan2019iotplatform}. \textbf{Consistent Data Synchronization:} Ensuring the dashboard always reflects the latest reading required consistent timestamp handling and backend validation \cite{thilakarathne2023fields}.


% ----------------------------------------------------------
% 4.5 Summary
% ----------------------------------------------------------
\subsection{Summary}
\label{sec:chapter4-summary}

This chapter presented the implementation details of the Smart AI 
Powered Agriculture System, including embedded firmware, backend 
services, mobile application logic, notifications, and AI-based 
crop recommendations \cite{ali2025soilml}. The described modules 
directly implement the requirements defined in Chapter~2 and 
conform to the system design presented in Chapter~3.

\setcounter{chapter}{4}